{\small
\begin{longtable}{@{}p{22mm}p{29mm}p{39mm}p{26mm}p{49mm}@{}}
\caption{Author-mapped GSE165816 DFU patient scores (7 healed, 4 not healed). Each score is the mean frozen fibroblast-associated topic loading over all retained cells in the listed foot specimens; pooling cells weights repeated specimens by retained cell count. Cells are counts, not independent outcome observations. Patient identifiers are study-derived pseudonyms and G labels identify specimens. The outcome is the source-study category within 12 weeks, not an individual healing day. Scores are on the topic-loading scale from zero to one.}\label{tab:p1patients}\\
\toprule
Patient & Outcome & Specimens & Cells & All-cell mean \\
\midrule\endfirsthead
\multicolumn{5}{l}{\tablename\ \thetable\ (continued)}\\
\toprule
Patient & Outcome & Specimens & Cells & All-cell mean \\
\midrule\endhead
\bottomrule\endfoot
S01 & not healed & G6 & 2279 & 0.0498 \\*
S02 & healed & G23, G4 & 7495 & 0.1496 \\
S03 & healed & G2 & 3951 & 0.0185 \\
S04 & healed & G15 & 2250 & 0.0758 \\
S05 & not healed & G9 & 4290 & 0.0346 \\
S06 & not healed & G33, G34 & 7681 & 0.0103 \\
S07 & not healed & G39 & 2481 & 0.0277 \\
S08 & healed & G7, G8 & 4192 & 0.0754 \\
S09 & healed & G49 & 2928 & 0.0534 \\
S10 & healed & G42 & 1838 & 0.0167 \\*
S11 & healed & G45 & 3374 & 0.0297 \\
\end{longtable}
}

{\small
\begin{longtable}{@{}p{25mm}p{20mm}p{49mm}p{37mm}p{34mm}@{}}
\caption{Historical scoring diagnostics across 9 healing and 5 non-healing GSE165816 specimens. AUC is the area under the ROC curve for pooled leave-one-sample-out scores; higher values favour healing. The percentile 95\% intervals use 4,000 arm-stratified resamples of fixed out-of-fold scores. One-sided permutation probabilities and null means rerun fold-dependent orientation and fitting. PCA is principal component analysis; NMF is nonnegative matrix factorization. Their raw coordinates are not calibrated across fitted folds, so these estimates cannot rank representations. The frozen topic encoder used discovery expression, and repeated patients limit independence. None of the six historical paired AUC-difference intervals excludes zero.}\label{tab:p1sampleauc}\\
\toprule
Readout & AUC & Bootstrap 95\% interval & Permutation p & Null mean \\
\midrule\endfirsthead
\multicolumn{5}{l}{\tablename\ \thetable\ (continued)}\\
\toprule
Readout & AUC & Bootstrap 95\% interval & Permutation p & Null mean \\
\midrule\endhead
\bottomrule\endfoot
Fibroblast topic & 0.822 & 0.556 to 1.000 & 0.0604 & 0.562 \\*
Module & 0.844 & 0.578 to 1.000 & 0.0424 & 0.497 \\
PCA & 0.533 & 0.156 to 0.889 & 0.4527 & 0.480 \\*
NMF & 0.878 & 0.655 to 1.000 & 0.0597 & 0.553 \\
\end{longtable}
}

{\small
\begin{longtable}{@{}p{45mm}p{25mm}p{59mm}p{38mm}@{}}
\caption{Historical patient-unit scoring diagnostics and paired anatomical sensitivity within GSE165816. The first four rows are pooled leave-one-patient-out AUCs for 7 healing and 4 non-healing patients, with 4,000 arm-stratified bootstrap draws of fixed scores and one-sided label-permutation probabilities. PCA and NMF denote principal component analysis and nonnegative matrix factorization; their coordinates are not comparable across fitted folds. These historical estimates are not uncertainty summaries for the revised held-pair procedure and do not rank representations. The final row is the foot-minus-forearm all-cell fibroblast-associated topic mean difference in 10 paired patients, with a 4,000-draw paired bootstrap interval and a two-sided exhaustive sign-flip probability using add-one correction. Its Benjamini--Hochberg q value across 15 topics is 0.0402. These row types have different estimands; neither supplies external validation.}\label{tab:p1patientauc}\\
\toprule
Readout & Estimate & Bootstrap 95\% interval & Permutation p \\
\midrule\endfirsthead
\multicolumn{4}{l}{\tablename\ \thetable\ (continued)}\\
\toprule
Readout & Estimate & Bootstrap 95\% interval & Permutation p \\
\midrule\endhead
\bottomrule\endfoot
Fibroblast topic & 0.7143 & 0.3571 to 1.0000 & 0.2991 \\*
Module & 0.7857 & 0.5000 to 1.0000 & 0.1662 \\
PCA & 0.5000 & 0.0000 to 1.0000 & 0.5025 \\
NMF & 0.1250 & 0.0000 to 0.4286 & 0.9751 \\*
foot minus forearm & 0.0476 & 0.0119 to 0.0916 & 0.0107 \\
\end{longtable}
}

\Needspace{155pt}
{\small
\begin{longtable}{@{}p{53mm}p{53mm}p{65mm}@{}}
\caption{Spearman correlations between the four representative sample-level prediction scores across the same 14 DFU specimens (9 healing, 5 non-healing). Fibroblast topic is the frozen all-cell loading, Module the specified fibro-inflammatory score, PCA principal component analysis and NMF nonnegative matrix factorization. Each row is one pair of readouts; values are descriptive point estimates without confidence intervals or multiplicity-adjusted tests. Related scores do not establish equivalent discrimination or rank methods.}\label{tab:p1correlations}\\
\toprule
Readout A & Readout B & Spearman correlation \\
\midrule\endfirsthead
\multicolumn{3}{l}{\tablename\ \thetable\ (continued)}\\
\toprule
Readout A & Readout B & Spearman correlation \\
\midrule\endhead
\bottomrule\endfoot
Fibroblast topic & Module & 0.7099 \\*
Fibroblast topic & PCA & 0.1736 \\
Fibroblast topic & NMF & 0.6067 \\
Module & PCA & 0.4022 \\
Module & NMF & 0.6756 \\*
PCA & NMF & 0.3200 \\
\end{longtable}
}

{\small
\begin{longtable}{@{}p{50mm}p{30mm}p{20mm}p{30mm}p{35mm}@{}}
\caption{Conditional next-cohort simulations for all-cell fibroblast-associated topic mean loading. Rows assume hypothetical independent observations, no missing outcomes, healing allocation near 9/14, and arm SDs 0.05170 and 0.01686 estimated from deterministic discovery specimens. These descriptive spreads include repeated specimens and are not patient-population variance estimates. Two-sided Welch tests assess candidate mean differences; Welch two-one-sided tests assess equivalence under a true difference of zero, at alpha=0.05. Total and arm columns give the simulated observation counts; power uses 4,000 draws per candidate size. Values near 0.80 have Monte Carlo error and are approximate search outputs. Candidate differences and margins are inputs, not clinical minimum important differences or validated recruitment targets.}\label{tab:p1design}\\
\toprule
Design & Difference / margin & Total & Healing / other & Estimated power \\
\midrule\endfirsthead
\multicolumn{5}{l}{\tablename\ \thetable\ (continued)}\\
\toprule
Design & Difference / margin & Total & Healing / other & Estimated power \\
\midrule\endhead
\bottomrule\endfoot
Detect difference & 0.02 & 100 & 64 / 36 & 0.797 \\*
Detect difference & 0.03 & 47 & 30 / 17 & 0.821 \\
Detect difference & 0.0457 & 21 & 14 / 7 & 0.811 \\
Detect difference & 0.06 & 14 & 9 / 5 & 0.804 \\
Detect difference & 0.08 & 9 & 6 / 3 & 0.807 \\
Detect difference & 0.1 & 8 & 5 / 3 & 0.880 \\
Equivalence under zero difference & 0.02 & 110 & 71 / 39 & 0.815 \\
Equivalence under zero difference & 0.03 & 49 & 32 / 17 & 0.826 \\
Equivalence under zero difference & 0.05 & 20 & 13 / 7 & 0.830 \\
Equivalence under zero difference & 0.075 & 11 & 7 / 4 & 0.880 \\*
Equivalence under zero difference & 0.1 & 8 & 5 / 3 & 0.915 \\
\end{longtable}
}

{\small
\begin{longtable}{@{}p{44mm}p{38mm}p{51mm}p{34mm}@{}}
\caption{Exploratory binormal planning simulation for paired AUC differences. Two candidate AUCs are centred on 0.80, separated by the stated gap, with within-class score-noise correlation 0.5045 and healing allocation near 9/14. Each searched total uses 4,000 simulation draws and a model-based normal approximation for the paired difference at alpha=0.05. Totals assume one observation per independent patient and no missing outcome. The gaps are hypothetical inputs; the approximate 80\% power search does not establish actual clinical performance, calibrated patient-test power or a validated recruitment target.}\label{tab:p1aucdesign}\\
\toprule
Candidate AUC gap & Total observations & Healing / non-healing & Estimated power \\
\midrule\endfirsthead
\multicolumn{4}{l}{\tablename\ \thetable\ (continued)}\\
\toprule
Candidate AUC gap & Total observations & Healing / non-healing & Estimated power \\
\midrule\endhead
\bottomrule\endfoot
0.05 & 706 & 454 / 252 & 0.800 \\*
0.1 & 181 & 116 / 65 & 0.814 \\
0.15 & 85 & 55 / 30 & 0.797 \\*
0.2 & 50 & 32 / 18 & 0.800 \\
\end{longtable}
}

{\small
\begin{longtable}{@{}p{63mm}p{28mm}p{39mm}p{37mm}@{}}
\caption{Association between mean fibroblast-associated topic loading within fibroblasts and the threshold-defined high-state fraction across 25 GSE165816 specimens. The Spearman row gives the observed correlation and percentile 95\% limits from 10,000 sample-row bootstrap draws. The leave-one-out row gives the minimum and maximum of 25 recalculations, not a confidence interval; a dash means no single point estimate applies. A separate 10,000-draw permutation analysis gave a two-sided add-one probability of 0.00010 and a central null interval from -0.4025 to 0.4072. Samples may share patients, and the two summaries share cells.}\label{tab:p1semantics}\\
\toprule
Summary & Estimate & Lower endpoint & Upper endpoint \\
\midrule\endfirsthead
\multicolumn{4}{l}{\tablename\ \thetable\ (continued)}\\
\toprule
Summary & Estimate & Lower endpoint & Upper endpoint \\
\midrule\endhead
\bottomrule\endfoot
spearman rho & 0.9325 & 0.7887 & 0.9828 \\*
leave one out rho & — & 0.9235 & 0.9566 \\
\end{longtable}
}

{\small
\begin{longtable}{@{}p{29mm}p{18mm}p{29mm}p{29mm}p{29mm}p{29mm}@{}}
\caption{External sample-level Spearman correlations with the frozen fibroblast-associated topic mean under exclusive lineage gates. Immune fraction divides immune-assigned cells by all retained cells; B/plasma share divides B/plasma cells by immune-assigned cells. Detected genes and unique molecular identifiers (UMIs) are separate per-sample medians over QC cells. All entries are descriptive point estimates without confidence intervals. Specimens are not verified independent patients, and no listed cohort supplies a qualifying patient-level healing outcome.}\label{tab:p1external}\\
\toprule
Cohort & Samples & Immune fraction & B/plasma share & Detected genes & UMIs \\
\midrule\endfirsthead
\multicolumn{6}{l}{\tablename\ \thetable\ (continued)}\\
\toprule
Cohort & Samples & Immune fraction & B/plasma share & Detected genes & UMIs \\
\midrule\endhead
\bottomrule\endfoot
GSE223964 & 8 & -0.024 & -0.071 & -0.429 & -0.476 \\*
GSE245703 & 12 & 0.846 & 0.077 & -0.427 & -0.308 \\*
GSE268834 & 8 & 0.714 & -0.500 & 0.238 & 0.262 \\
\end{longtable}
}

{\small
\begin{longtable}{@{}p{76mm}p{20mm}p{25mm}p{22mm}p{22mm}@{}}
\caption{Complete exploratory covariate screen in 25 discovery specimens. Each row correlates the fibroblast high-state fraction, using the historical rule loading at least 0.184, with one of 31 covariates. Lineage fractions use all retained cells; technical summaries and topic means use fibroblasts. Values are Spearman rho, two-sided asymptotic p and Benjamini--Hochberg q across all 31 tests, ordered by absolute rho; 14 rows have q below 0.05. UMI means unique molecular identifier. The G label is an identifier, not an established batch or time variable. Repeated patients invalidate the naive independent-row premise of these historical p values; their BH transforms do not establish false-discovery-rate control or 14 independent discoveries. Shared cells and compositional constraints further limit interpretation; these are descriptive row-level summaries, not causal or healing-validation tests.}\label{tab:p1covariates}\\
\toprule
Covariate & Specimens & Spearman rho & Raw p & BH q \\
\midrule\endfirsthead
\multicolumn{5}{l}{\tablename\ \thetable\ (continued)}\\
\toprule
Covariate & Specimens & Spearman rho & Raw p & BH q \\
\midrule\endhead
\bottomrule\endfoot
B/plasma fraction & 25 & 0.8279 & 3.26e-07 & 1.01e-05 \\*
Fibroblast CFD/APOD mean & 25 & -0.7664 & 7.95e-06 & 0.000123 \\
Fibroblast PTPRC-positive fraction & 25 & 0.7267 & 3.88e-05 & 0.000401 \\
Fibroblast median UMI count & 25 & 0.6632 & 0.000302 & 0.00234 \\
Fibroblast IGFBP7/TAGLN mean & 25 & 0.6514 & 0.000421 & 0.00261 \\
Sweat-gland fraction & 25 & -0.5832 & 0.00222 & 0.0114 \\
Fibroblast FOS/IGFBP7 mean & 25 & -0.5687 & 0.00302 & 0.0134 \\
Fibroblast DCN/CFD mean & 25 & -0.5592 & 0.00366 & 0.0142 \\
Fibroblast IGFBP7/IGFBP5 mean & 25 & 0.5316 & 0.00624 & 0.0215 \\
Fibroblast COL1A1/COL1A2 mean & 25 & 0.5238 & 0.00721 & 0.0223 \\
Keratinocyte fraction & 25 & -0.4836 & 0.0143 & 0.0404 \\
Fibroblast HLA-DRA/CD74 mean & 25 & 0.4726 & 0.0171 & 0.0441 \\
T-cell fraction & 25 & 0.4647 & 0.0193 & 0.0459 \\
Fibroblast TAGLN/ACTA2 mean & 25 & 0.4552 & 0.0222 & 0.0492 \\
Melanocyte fraction & 25 & -0.4442 & 0.0261 & 0.054 \\
Fibroblast JUNB/DNAJB1 mean & 25 & 0.4182 & 0.0375 & 0.0726 \\
Mast-cell fraction & 25 & -0.3977 & 0.0489 & 0.0893 \\
Myeloid fraction & 25 & 0.3891 & 0.0546 & 0.094 \\
Numeric G specimen label & 25 & -0.3607 & 0.0765 & 0.125 \\
Endothelial fraction & 25 & -0.2347 & 0.259 & 0.401 \\
Total retained cells & 25 & -0.2276 & 0.274 & 0.404 \\
Fibroblast DNAJB1/FOS mean & 25 & 0.1977 & 0.344 & 0.484 \\
Fibroblast KRT14/KRT10 mean & 25 & 0.1457 & 0.487 & 0.63 \\
Fibroblast dissociation score & 25 & 0.1449 & 0.489 & 0.63 \\
Lymphatic endothelial fraction & 25 & -0.1331 & 0.526 & 0.63 \\
Fibroblast fraction & 25 & 0.1292 & 0.538 & 0.63 \\
Fibroblast FOS/JUNB mean & 25 & -0.1260 & 0.548 & 0.63 \\
Retained fibroblasts & 25 & 0.0851 & 0.686 & 0.759 \\
Pericyte/smooth-muscle fraction & 25 & -0.0473 & 0.823 & 0.879 \\
Fibroblast IGFBP7/SPARCL1 mean & 25 & -0.0354 & 0.866 & 0.895 \\*
Fibroblast DCD/SAT1 mean & 25 & -0.0236 & 0.911 & 0.911 \\
\end{longtable}
}

{\small
\begin{longtable}{@{}p{64mm}p{13mm}p{22mm}p{23mm}p{23mm}p{20mm}@{}}
\caption{Immune-transcript exclusion and residual-mixture controls. BIC denotes Bayesian information criterion, fitted to sample high-state fractions or their regression residuals; the difference is BIC for one Gaussian minus BIC for two. Positive differences favour two components. The first two rows share 24 matched specimens and have fraction correlation 0.9887. Residual regressions use all 25 specimens: immune covariates are B/plasma, T-cell and myeloid fractions plus immune-transcript positivity; placebos are pericyte/smooth-muscle, endothelial, melanocyte and sweat-gland fractions. R squared is variance explained by each regression and is inapplicable to unadjusted rows (dash). Both residual differences are negative, so the analysis does not establish immune-specific separation.}\label{tab:p1ambient}\\
\toprule
Summary & n & R squared & BIC one & BIC two & Difference \\
\midrule\endfirsthead
\multicolumn{6}{l}{\tablename\ \thetable\ (continued)}\\
\toprule
Summary & n & R squared & BIC one & BIC two & Difference \\
\midrule\endhead
\bottomrule\endfoot
Before immune-transcript exclusion & 24 & — & 12.417 & -28.569 & 40.986 \\*
After immune-transcript exclusion & 24 & — & 5.342 & -51.341 & 56.683 \\
Immune-covariate residuals & 25 & 0.4953 & -3.900 & 0.031 & -3.931 \\*
Placebo-covariate residuals & 25 & 0.3549 & 2.235 & 10.433 & -8.198 \\
\end{longtable}
}

{\small
\begin{longtable}{@{}p{29mm}p{49mm}p{24mm}p{42mm}p{21mm}@{}}
\caption{Historical sample-level lineage and loading contrasts. Ratios divide the healing-arm mean by the non-healing-arm mean; intervals are percentile 95\% limits from 20,000 within-arm sample bootstraps, and p values are two-sided Mann--Whitney tests. GSE165816 uses 9 versus 5 DFU specimens; GSE231643 uses 5 versus 3 title-labelled specimens. The discovery ratios come from the saved lineage-projection analysis and are distinct from the later deterministic all-cell patient contrast. Fibroblast fraction uses all cells as denominator; fibroblast mean uses only assigned fibroblasts. These are separate estimands. Repeated discovery patients and unavailable external patient mappings prevent independent outcome validation.}\label{tab:p1lineage}\\
\toprule
Cohort & Readout & Mean ratio & Bootstrap 95\% interval & Raw p \\
\midrule\endfirsthead
\multicolumn{5}{l}{\tablename\ \thetable\ (continued)}\\
\toprule
Cohort & Readout & Mean ratio & Bootstrap 95\% interval & Raw p \\
\midrule\endhead
\bottomrule\endfoot
GSE165816 & Fibroblast fraction & 1.425 & 0.873 to 2.298 & 0.438 \\*
GSE165816 & All-cell topic mean & 2.724 & 1.369 to 5.732 & 0.0599 \\
GSE165816 & Within-fibroblast topic mean & 2.889 & 1.059 to 13.966 & 0.042 \\
GSE231643 & Fibroblast fraction & 2.051 & 1.128 to 4.184 & 0.0714 \\
GSE231643 & All-cell topic mean & 2.564 & 1.310 to 8.079 & 0.143 \\*
GSE231643 & Within-fibroblast topic mean & 1.660 & 1.007 to 6.102 & 0.143 \\
\end{longtable}
}

{\small
\begin{longtable}{@{}p{31mm}p{59mm}p{26mm}p{51mm}@{}}
\caption{Historical paired AUC differences for the pooled-fold scoring procedure at each analysis unit. Each contrast subtracts the second leave-one-out readout from the first. Percentile 95\% intervals use 4,000 common-index, arm-stratified bootstrap draws of fixed out-of-fold scores: 9 versus 5 specimens or 7 versus 4 mapped patients. PCA and NMF denote principal component analysis and nonnegative matrix factorization. No sample-unit interval excludes zero; patient-unit Fibroblast topic minus NMF and Module minus NMF intervals do. Because raw PCA/NMF coordinates can change between fitted folds, neither exclusion from zero supports representation superiority. These historical intervals are unadjusted, omit refitting uncertainty and do not describe the revised within-model held-pair estimand.}\label{tab:p1aucdifferences}\\
\toprule
Unit & AUC contrast & Difference & Bootstrap 95\% interval \\
\midrule\endfirsthead
\multicolumn{4}{l}{\tablename\ \thetable\ (continued)}\\
\toprule
Unit & AUC contrast & Difference & Bootstrap 95\% interval \\
\midrule\endhead
\bottomrule\endfoot
14 specimens & Fibroblast topic minus Module & -0.0222 & -0.2889 to 0.2222 \\*
14 specimens & Fibroblast topic minus PCA & 0.2889 & -0.0889 to 0.6667 \\
14 specimens & Fibroblast topic minus NMF & -0.0556 & -0.3556 to 0.2222 \\
14 specimens & Module minus PCA & 0.3111 & 0.0000 to 0.6667 \\
14 specimens & Module minus NMF & -0.0333 & -0.2889 to 0.2000 \\
14 specimens & PCA minus NMF & -0.3444 & -0.7111 to 0.0000 \\
11 patients & Fibroblast topic minus Module & -0.0714 & -0.4286 to 0.2857 \\
11 patients & Fibroblast topic minus PCA & 0.2143 & -0.4643 to 0.9286 \\
11 patients & Fibroblast topic minus NMF & 0.5893 & 0.1429 to 0.9643 \\
11 patients & Module minus PCA & 0.2857 & -0.1786 to 0.7857 \\
11 patients & Module minus NMF & 0.6607 & 0.2857 to 0.9652 \\*
11 patients & PCA minus NMF & 0.3750 & 0.0000 to 0.7857 \\
\end{longtable}
}

{\small
\begin{longtable}{@{}p{42mm}p{35mm}p{48mm}p{21mm}p{19mm}@{}}
\caption{Complete paired anatomical control across 15 topics in 10 author-mapped patients. Differences are mean all-cell loadings in foot minus forearm, after cell-count-weighted collapse of repeated foot specimens. Intervals are percentile 95\% limits from 4,000 paired-patient bootstrap draws. Two-sided p values enumerate all 1,024 sign allocations and retain add-one correction; q values use Benjamini--Hochberg across 15 topics. The TIMP1/COL3A1, COL1A1/COL1A2, CFD/APOD and IGFBP7/TAGLN loadings have q below 0.05. Labels list the two leading decoder genes; they do not assert cell-type identity. Bootstrap intervals and permutation tests need not be dual. This is an anatomical sensitivity control within the discovery source, not a healing endpoint.}\label{tab:p1anatomy}\\
\toprule
Leading genes & Foot minus forearm & Bootstrap 95\% interval & Sign-flip p & BH q \\
\midrule\endfirsthead
\multicolumn{5}{l}{\tablename\ \thetable\ (continued)}\\
\toprule
Leading genes & Foot minus forearm & Bootstrap 95\% interval & Sign-flip p & BH q \\
\midrule\endhead
\bottomrule\endfoot
TIMP1/COL3A1 & 0.0476 & 0.0119 to 0.0916 & 0.0107 & 0.0402 \\*
JUNB/DNAJB1 & 0.0103 & -0.0616 to 0.0793 & 0.799 & 0.956 \\
COL1A1/COL1A2 & 0.0914 & 0.0535 to 0.1376 & 0.00293 & 0.022 \\
IGFBP7/IGFBP5 & 0.0188 & 0.0041 to 0.0361 & 0.0634 & 0.119 \\
TAGLN/ACTA2 & 0.0016 & -0.0200 to 0.0241 & 0.893 & 0.956 \\
KRT14/KRT10 & -0.2491 & -0.4274 to -0.0682 & 0.0361 & 0.0774 \\
CFD/APOD & -0.0362 & -0.0567 to -0.0169 & 0.0107 & 0.0402 \\
DCN/CFD & -0.0159 & -0.0383 to 0.0042 & 0.208 & 0.312 \\
DCD/SAT1 & -0.0001 & -0.0083 to 0.0098 & 0.975 & 0.975 \\
IGFBP7/SPARCL1 & 0.0398 & -0.0015 to 0.0897 & 0.134 & 0.223 \\
FOS/IGFBP7 & -0.0038 & -0.0087 to 0.0025 & 0.245 & 0.334 \\
HLA-DRA/CD74 & 0.0111 & -0.0591 to 0.0984 & 0.854 & 0.956 \\
IGFBP7/TAGLN & 0.0715 & 0.0436 to 0.1039 & 0.00293 & 0.022 \\
FOS/JUNB & 0.0068 & 0.0025 to 0.0118 & 0.0185 & 0.0556 \\*
DNAJB1/FOS & 0.0061 & 0.0020 to 0.0108 & 0.0224 & 0.0561 \\
\end{longtable}
}

{\small
\begin{longtable}{@{}p{33mm}p{23mm}p{33mm}p{33mm}p{43mm}@{}}
\caption{Complete Welch two-one-sided-test sensitivity grids for the all-cell mean fibroblast-associated topic contrast. The symmetric margin is on the loading scale, and both lower and upper one-sided p values must be below 0.05 for the final column to be Yes. Sample rows use 9 versus 5 specimens; patient rows use 7 versus 4 mapped patients with cell-count-weighted repeat collapse. Margins were examined retrospectively and have no externally justified clinical meaning. A Yes entry therefore records a mathematical sensitivity result, not demonstrated clinical equivalence.}\label{tab:p1equivalence}\\
\toprule
Unit & Margin & Lower-test p & Upper-test p & Both below 0.05 \\
\midrule\endfirsthead
\multicolumn{5}{l}{\tablename\ \thetable\ (continued)}\\
\toprule
Unit & Margin & Lower-test p & Upper-test p & Both below 0.05 \\
\midrule\endhead
\bottomrule\endfoot
14 specimens & 0.05 & 0.000198 & 0.413 & No \\*
14 specimens & 0.10 & 5.59e-06 & 0.0077 & Yes \\
14 specimens & 0.20 & 3.66e-08 & 3.31e-06 & Yes \\
14 specimens & 0.30 & 1.14e-09 & 2.6e-08 & Yes \\
14 specimens & 0.50 & 9.95e-12 & 6.75e-11 & Yes \\
11 patients & 0.05 & 0.00172 & 0.159 & No \\
11 patients & 0.10 & 7.56e-05 & 0.00323 & Yes \\
11 patients & 0.20 & 1.09e-06 & 1.02e-05 & Yes \\
11 patients & 0.30 & 6.42e-08 & 3e-07 & Yes \\*
11 patients & 0.50 & 1.45e-09 & 3.72e-09 & Yes \\
\end{longtable}
}

{\small
\begin{longtable}{@{}p{43mm}p{24mm}p{21mm}p{39mm}p{39mm}@{}}
\caption{Patient-level robustness in the discovery cohort. The first row correlates fibroblast mean loading with the high-state fraction after weighting repeated specimen summaries by fibroblast count within 20 patients (25 specimens). Its percentile 95\% interval uses 10,000 patient resamples, seed 217. The second row uses the previously defined all-cell healing contrast in 7 healing and 4 non-healing patients; its original interval remains in the outcome table. Omission ranges give the minimum and maximum after deleting one patient at a time (20 or 11 omissions), not confidence limits. Neither analysis refits the expression model, changes the state cut, adds patients, or provides external validation.}\label{tab:p1robustness}\\
\toprule
Quantity & Unit & Estimate & Patient bootstrap 95\% & Patient-omission range \\
\midrule\endfirsthead
\multicolumn{5}{l}{\tablename\ \thetable\ (continued)}\\
\toprule
Quantity & Unit & Estimate & Patient bootstrap 95\% & Patient-omission range \\
\midrule\endhead
\bottomrule\endfoot
Mean versus high-state fraction & 20 patients & 0.9372 & 0.7760 to 0.9903 & 0.9266 to 0.9643 \\*
Healing minus non-healing & 11 patients & 0.0293 & Not recalculated & 0.0143 to 0.0365 \\
\end{longtable}
}

{\small
\begin{longtable}{@{}p{12mm}p{46mm}p{25mm}p{26mm}p{21mm}p{35mm}@{}}
\caption{Topic-capacity sensitivity after tracking the TIMP1-carrying topic at each K. Gene-based descriptions identify the tracked loading; each fit supplies its own two-Gaussian midpoint cut. Semantic rho uses 25 specimens; the unadjusted healing-arm p value uses 9 versus 5 high-state-fraction summaries. The final column asks whether the maximum fraction occurs in healthy skin. At K=10 the tracked topic is a different IGFBP-rich programme and the counterexample fails. Topics across capacities are not assumed to have identical biological meaning.}\label{tab:p1capacity}\\
\toprule
K & Programme description & Derived cut & Semantic rho & Healing p & Healthy maximum \\
\midrule\endfirsthead
\multicolumn{6}{l}{\tablename\ \thetable\ (continued)}\\
\toprule
K & Programme description & Derived cut & Semantic rho & Healing p & Healthy maximum \\
\midrule\endhead
\bottomrule\endfoot
10 & IGFBP-rich & 0.1454 & 0.9469 & 0.147 & No \\*
15 & CHI3L1/COL3A1-rich & 0.1691 & 0.9316 & 0.109 & Yes \\*
20 & CHI3L1/COL3A1-rich & 0.1695 & 0.9236 & 0.109 & Yes \\
\end{longtable}
}

{\small
\begin{longtable}{@{}p{27mm}p{31mm}p{43mm}p{25mm}p{39mm}@{}}
\caption{All 18 state-cut sensitivity evaluations for the frozen K=15 topic. Spearman rho correlates fibroblast mean loading with the threshold-defined high-state fraction across 25 specimens. The ratio compares mean high-state fractions in 9 healing versus 5 non-healing specimens; p is the unadjusted two-sided Mann--Whitney probability. The last column gives the group containing the largest fraction. The grid includes the fitted cut rounded to 0.1843. Four cuts have p below 0.05; these are dependent sensitivity analyses, not 18 independent confirmations.}\label{tab:p1cuts}\\
\toprule
State cut & Semantic rho & Healing / non-healing & Raw p & Largest fraction \\
\midrule\endfirsthead
\multicolumn{5}{l}{\tablename\ \thetable\ (continued)}\\
\toprule
State cut & Semantic rho & Healing / non-healing & Raw p & Largest fraction \\
\midrule\endhead
\bottomrule\endfoot
0.0800 & 0.9326 & 2.966 & 0.042 & Healthy \\*
0.1000 & 0.9414 & 2.959 & 0.042 & Healthy \\
0.1200 & 0.9091 & 3.011 & 0.042 & Healthy \\
0.1400 & 0.9301 & 2.992 & 0.042 & Healthy \\
0.1600 & 0.9294 & 2.981 & 0.0599 & Healthy \\
0.1800 & 0.9349 & 3.001 & 0.0599 & Healthy \\
0.1843 & 0.9325 & 2.975 & 0.0599 & Healthy \\
0.2000 & 0.9349 & 2.987 & 0.0599 & Healthy \\
0.2200 & 0.9341 & 3.124 & 0.0599 & Healthy \\
0.2400 & 0.9228 & 3.352 & 0.0617 & Healthy \\
0.2600 & 0.9228 & 3.424 & 0.0617 & Healthy \\
0.2800 & 0.9236 & 3.465 & 0.0617 & Healthy \\
0.3000 & 0.9236 & 3.499 & 0.0617 & Healthy \\
0.3200 & 0.9479 & 3.741 & 0.0617 & Healthy \\
0.3400 & 0.9463 & 3.699 & 0.0617 & Healthy \\
0.3600 & 0.9401 & 4.040 & 0.0608 & Healthy \\
0.3800 & 0.9182 & 3.945 & 0.0608 & Healthy \\*
0.4000 & 0.9182 & 3.889 & 0.0608 & Healthy \\
\end{longtable}
}

{\small
\begin{longtable}{@{}p{63mm}p{55mm}p{49mm}@{}}
\caption{Exact mean decomposition on the saved 19,410 discovery fibroblasts in 25 specimens. The cut is 0.184295848; empirical bin means are distinct from fitted Gaussian-component means. The constant-state approximation uses the cell-pooled low/high means, and its residual is the observed specimen mean minus that approximation. RMS denotes root-mean-square residual with equal specimen weights. Ranges exclude undefined means for empty bins; zero-weight terms contribute zero to the identity. Both exact identities have maximum numerical discrepancy below $10^{-15}$. Values are on the zero-to-one loading scale, apart from the count row. They describe the frozen sample, without confidence intervals, a test of invariant state intensity or causal variance attribution.}\label{tab:p1math}\\
\toprule
Quantity & Value or range & Denominator \\
\midrule\endfirsthead
\multicolumn{3}{l}{\tablename\ \thetable\ (continued)}\\
\toprule
Quantity & Value or range & Denominator \\
\midrule\endhead
\bottomrule\endfoot
Cell-pooled low-bin mean & 0.02380 & All 19,410 fibroblasts \\*
Cell-pooled high-bin mean & 0.51301 & All 19,410 fibroblasts \\
Specimen low-bin means & 0.01204 to 0.10213 & 25 nonempty low bins \\
Specimen high-bin means & 0.22645 to 0.57777 & 16 nonempty high bins \\
Specimens with an empty bin & 9 & 9 of 25; high bin empty \\
Within-state residual range & -0.02434 to 0.04377 & All 25 specimens \\*
Within-state residual RMS & 0.01575 & Equal specimen weights \\
\end{longtable}
}

{\small
\begin{longtable}{@{}p{100mm}p{67mm}@{}}
\caption{Independent arithmetic check of the all-cell mean-loading contrast in 11 patients (7 healing, 4 non-healing). Exhaustive allocation includes the observed labels and numerical ties: 90/330 allocations are at least as extreme in absolute mean difference, whereas the retained add-one convention gives 91/331. These probabilities assume observational label exchangeability and do not identify a causal effect. Welch quantities use sample variances within the two outcome groups. Standard error, interval and margin are on the loading scale; degrees of freedom and probabilities are dimensionless. A margin strictly greater than the infimum would be needed to reject both one-sided tests at 0.05. The data-derived boundary is not a clinical equivalence margin.}\label{tab:p1exact}\\
\toprule
Quantity & Value \\
\midrule\endfirsthead
\multicolumn{2}{l}{\tablename\ \thetable\ (continued)}\\
\toprule
Quantity & Value \\
\midrule\endhead
\bottomrule\endfoot
All label allocations & 330 \\*
At least as extreme as observed & 90 \\
Exhaustive two-sided probability & 0.272727 \\
Recorded add-one probability & 0.274924 \\
Welch standard error & 0.019446 \\
Welch degrees of freedom & 8.113833 \\
Two-sided 90\% Welch interval & -0.006816 to 0.065375 \\*
Equivalence-margin infimum & 0.065375 \\
\end{longtable}
}

{\small
\begin{longtable}{@{}p{15mm}p{25mm}p{18mm}p{23mm}p{27mm}p{25mm}p{27mm}@{}}
\caption{Observed-cell extension: mapped patient-level tissue measurements from 25 selected foot specimens. Patient labels are stable pseudonyms within this extension. The 20 rows include 9 non-diabetic participants, who do not enter healing tests, and 11 DFU patients (7 healed and 4 not healed). Repeated specimens are pooled by fibroblast count for the mean and high-state fraction, and by all-cell count for fibroblast composition. High state uses the unchanged discovery midpoint. Specimens and fibroblasts are counts; the three readouts are unitless fractions or loadings on zero-to-one scales.}\label{tab:p1biologicalpatients}\\
\toprule
Patient & Group & Specimens & Fibroblasts & Fibroblast mean & High fraction & Fibroblast fraction \\
\midrule\endfirsthead
\multicolumn{7}{l}{\tablename\ \thetable\ (continued)}\\
\toprule
Patient & Group & Specimens & Fibroblasts & Fibroblast mean & High fraction & Fibroblast fraction \\
\midrule\endhead
\bottomrule\endfoot
01 & Not healed & 1 & 221 & 0.2496 & 0.5113 & 0.0970 \\*
02 & Healed & 2 & 2751 & 0.3473 & 0.5863 & 0.3670 \\
03 & Healed & 1 & 1238 & 0.0274 & 0.0226 & 0.3133 \\
04 & Healed & 1 & 356 & 0.3521 & 0.6994 & 0.1582 \\
05 & Not healed & 1 & 843 & 0.0125 & 0.0000 & 0.1965 \\
06 & Not healed & 2 & 1552 & 0.0179 & 0.0032 & 0.2021 \\
07 & Not healed & 1 & 872 & 0.0547 & 0.0837 & 0.3515 \\
08 & Healed & 2 & 1605 & 0.1648 & 0.2960 & 0.3829 \\
09 & Healed & 1 & 363 & 0.2087 & 0.3719 & 0.1240 \\
10 & Healed & 1 & 170 & 0.0705 & 0.0824 & 0.0925 \\
11 & Healed & 1 & 1821 & 0.0393 & 0.0483 & 0.5397 \\
12 & Non-diabetic & 1 & 1020 & 0.4663 & 0.8412 & 0.3491 \\
13 & Non-diabetic & 1 & 215 & 0.0295 & 0.0186 & 0.0866 \\
14 & Non-diabetic & 2 & 1921 & 0.0137 & 0.0000 & 0.2509 \\
15 & Non-diabetic & 1 & 1152 & 0.0120 & 0.0000 & 0.3904 \\
16 & Non-diabetic & 2 & 1562 & 0.0153 & 0.0000 & 0.1873 \\
17 & Non-diabetic & 1 & 153 & 0.0136 & 0.0000 & 0.1939 \\
18 & Non-diabetic & 1 & 974 & 0.0131 & 0.0010 & 0.1804 \\
19 & Non-diabetic & 1 & 425 & 0.0135 & 0.0000 & 0.2941 \\*
20 & Non-diabetic & 1 & 196 & 0.0205 & 0.0000 & 0.1142 \\
\end{longtable}
}

{\small
\begin{longtable}{@{}p{34mm}p{34mm}p{23mm}p{36mm}p{17mm}p{17mm}@{}}
\caption{Patient-unit extension of the three fibroblast readouts. Differences are healed minus not healed in the same 7 and 4 DFU patients. Percentile intervals use 10,000 within-arm patient resamples, seed 17. Exact two-sided probabilities enumerate all 330 label allocations including ties, without add-one correction; q values apply Benjamini--Hochberg adjustment to the three readouts within each weighting variant. These are retrospective observational comparisons without a clinical equivalence margin. Every DFU specimen already has at least 100 fibroblasts, so the 100-cell eligibility variant reproduces the pooled-cell comparison exactly and duplicate rows are omitted. Repeated-cell and specimen weighting affect only aggregation within a patient.}\label{tab:p1biologicalcontrasts}\\
\toprule
Aggregation & Readout & Difference & Bootstrap 95\% interval & Exact p & BH q \\
\midrule\endfirsthead
\multicolumn{6}{l}{\tablename\ \thetable\ (continued)}\\
\toprule
Aggregation & Readout & Difference & Bootstrap 95\% interval & Exact p & BH q \\
\midrule\endhead
\bottomrule\endfoot
Pooled cells & Fibroblast mean & 0.0892 & -0.0486 to 0.2177 & 0.3121 & 0.4697 \\*
Pooled cells & High-state fraction & 0.1514 & -0.1331 to 0.4153 & 0.3697 & 0.4697 \\
Pooled cells & Fibroblast composition & 0.0708 & -0.0708 to 0.2146 & 0.4697 & 0.4697 \\
Equal specimens & Fibroblast mean & 0.0974 & -0.0421 to 0.2260 & 0.2788 & 0.4848 \\
Equal specimens & High-state fraction & 0.1683 & -0.1170 to 0.4310 & 0.3424 & 0.4848 \\*
Equal specimens & Fibroblast composition & 0.0679 & -0.0724 to 0.2114 & 0.4848 & 0.4848 \\
\end{longtable}
}

{\small
\begin{longtable}{@{}p{28mm}p{24mm}p{32mm}p{33mm}p{47mm}@{}}
\caption{Patient-level threshold sensitivity in the same 7 healed and 4 non-healed DFU patients. Each patient has equal group weight; recovered fibroblasts are pooled within patient. Differences are healed minus non-healed high-state fractions. Exact two-sided p values enumerate all 330 patient-label allocations, including ties, without add-one correction. The support-wide p compares each absolute contrast with the maximum absolute patient-group CDF difference over every distinct observed loading in each allocation; it accounts for searching across that complete support under the exchangeability null. The threshold-free maximum is 0.3045 with exact p=0.1848. Original denotes the unchanged fitted midpoint. These retrospective diagnostics are not independent validation.}\label{tab:p1thresholdpatient}\\
\toprule
Threshold & Original & Difference & Exact p & Support-wide p \\
\midrule\endfirsthead
\multicolumn{5}{l}{\tablename\ \thetable\ (continued)}\\
\toprule
Threshold & Original & Difference & Exact p & Support-wide p \\
\midrule\endhead
\bottomrule\endfoot
0.0500 & False & 0.2485 & 0.2242 & 0.3182 \\*
0.1000 & False & 0.1958 & 0.3030 & 0.4515 \\
0.1500 & False & 0.1682 & 0.3424 & 0.5606 \\
0.1843 & True & 0.1514 & 0.3697 & 0.6152 \\
0.2000 & False & 0.1445 & 0.3727 & 0.6424 \\
0.3000 & False & 0.1300 & 0.3727 & 0.7182 \\
0.4000 & False & 0.1029 & 0.3909 & 0.8061 \\
0.5000 & False & 0.0819 & 0.3667 & 0.8758 \\*
0.6000 & False & 0.0570 & 0.3091 & 0.9485 \\
\end{longtable}
}

{\small
\begin{longtable}{@{}p{30mm}p{31mm}p{23mm}p{32mm}p{32mm}p{16mm}@{}}
\caption{Cell-count and repeated-specimen sensitivity at the fixed original threshold. Equal-cell rows use 250 draws of 100 fibroblasts without replacement from each DFU specimen, seed 29, followed by equal specimen weighting within each patient. The other rows enumerate every one-specimen-per-patient choice across the three repeated patients. Every comparison retains 7 healed and 4 non-healed patients and recomputes the exact 330-allocation probability. Difference ranges are healed minus non-healed; positive counts cover the listed number of comparisons. Neither range is a confidence interval. Composition is omitted from equal-fibroblast sampling because that selected subset no longer estimates the original all-cell denominator.}\label{tab:p1samplingsensitivity}\\
\toprule
Sensitivity & Readout & Comparisons & Difference range & Exact p range & Positive \\
\midrule\endfirsthead
\multicolumn{6}{l}{\tablename\ \thetable\ (continued)}\\
\toprule
Sensitivity & Readout & Comparisons & Difference range & Exact p range & Positive \\
\midrule\endhead
\bottomrule\endfoot
100 cells per specimen & Fibroblast mean & 250 & 0.0749 to 0.1233 & 0.1273 to 0.3879 & 250/250 \\*
100 cells per specimen & High-state fraction & 250 & 0.1098 to 0.2229 & 0.1515 to 0.5667 & 250/250 \\
One specimen per patient & Fibroblast mean & 8 & 0.0777 to 0.1170 & 0.2333 to 0.3303 & 8/8 \\
One specimen per patient & High-state fraction & 8 & 0.1284 to 0.2081 & 0.2667 to 0.4394 & 8/8 \\*
One specimen per patient & Fibroblast composition & 8 & 0.0440 to 0.0918 & 0.3909 to 0.6515 & 8/8 \\
\end{longtable}
}

{\small
\begin{longtable}{@{}p{73mm}p{47mm}p{47mm}@{}}
\caption{Raw-count column-shuffle control on 6,000 cells, with 240 cells sampled from each of 25 specimens. Every one of 23,572 gene-value multisets is preserved exactly before normalization. Both branches use the same frozen encoder and deterministic projection. HHI is the sum of squared topic proportions; lower HHI and higher perplexity describe reduced concentration. Library sizes and detected-gene counts use all source genes. Detection fraction divides the detected-gene count by 23,572. Means are preserved by column marginals, but row-wise variation narrows. These descriptive values neither isolate biological co-expression from technical row structure nor show rejection by the historical integrity gate.}\label{tab:p1rawcontrol}\\
\toprule
Quantity & Original counts & Shuffled counts \\
\midrule\endfirsthead
\multicolumn{3}{l}{\tablename\ \thetable\ (continued)}\\
\toprule
Quantity & Original counts & Shuffled counts \\
\midrule\endhead
\bottomrule\endfoot
Mean topic perplexity & 3.4383 & 11.2306 \\*
Mean HHI & 0.5740 & 0.1467 \\
Library size, mean & 9275.2065 & 9275.2065 \\
Library size, SD & 10907.9970 & 1500.9520 \\
Detected genes, mean & 2312.1958 & 2312.1958 \\
Detected genes, SD & 1367.1079 & 38.9770 \\
Detection fraction, mean & 0.0981 & 0.0981 \\*
Detection fraction, SD & 0.0580 & 0.0017 \\
\end{longtable}
}

{\small
\begin{longtable}{@{}p{43mm}p{51mm}p{25mm}p{24mm}p{24mm}@{}}
\caption{Marker-reference sensitivity in 11 DFU patients. Re-estimating marker means and SDs within each specimen changes 7,044 of 76,461 cell labels (9.21\%); fibroblast calls change from 19,410 to 19,711, with 1,119 lost and 1,420 gained. The maximum specimen changes are 0.0371 for the fibroblast mean and 0.0695 for high-state fraction. Listed differences are healed minus non-healed after count-weighted patient aggregation; exact probabilities enumerate 330 allocations, and BH q adjusts the three readouts within each gate. Applying the frozen pooled reference per specimen reproduces all original labels and is omitted as a duplicate implementation check. No independent cell-identity truth set selects the correct gate.}\label{tab:p1gatecontrol}\\
\toprule
Marker reference & Readout & Difference & Exact p & BH q \\
\midrule\endfirsthead
\multicolumn{5}{l}{\tablename\ \thetable\ (continued)}\\
\toprule
Marker reference & Readout & Difference & Exact p & BH q \\
\midrule\endhead
\bottomrule\endfoot
Pooled reference & Fibroblast mean & 0.0892 & 0.3121 & 0.4697 \\*
Pooled reference & High-state fraction & 0.1514 & 0.3697 & 0.4697 \\
Pooled reference & Fibroblast composition & 0.0708 & 0.4697 & 0.4697 \\
Per-specimen reference & Fibroblast mean & 0.0887 & 0.3182 & 0.4970 \\
Per-specimen reference & High-state fraction & 0.1558 & 0.3667 & 0.4970 \\*
Per-specimen reference & Fibroblast composition & 0.0553 & 0.4970 & 0.4970 \\
\end{longtable}
}

{\small
\begin{longtable}{@{}p{54mm}p{28mm}p{54mm}p{31mm}@{}}
\caption{Secondary DFU-only specimen-label calibration. All 11 healthy specimens are excluded before enumerating the 2,002 allocations of 9 healed and 5 non-healed specimen labels. Differences are equal-specimen arm means, healed minus non-healed. The central 95\% null interval describes the allocation distribution and is not an effect confidence interval. Exact two-sided probabilities include the observed allocation and ties without add-one correction. Repeated specimens share patients, so the 11-patient comparisons remain primary.}\label{tab:p1dfunull}\\
\toprule
Readout & Difference & Central null 95\% & Exact p \\
\midrule\endfirsthead
\multicolumn{4}{l}{\tablename\ \thetable\ (continued)}\\
\toprule
Readout & Difference & Central null 95\% & Exact p \\
\midrule\endhead
\bottomrule\endfoot
Fibroblast mean & 0.1335 & -0.1652 to 0.1623 & 0.0894 \\*
High-state fraction & 0.2380 & -0.3167 to 0.3222 & 0.1199 \\*
Fibroblast composition & 0.0888 & -0.1540 to 0.1460 & 0.2697 \\
\end{longtable}
}

{\small
\begin{longtable}{@{}p{24mm}p{25mm}p{25mm}p{35mm}p{29mm}p{23mm}@{}}
\caption{Disagreement between five repeated-specimen pairs from five patients, including three DFU patients. G labels retain source specimen identifiers. Mean and high-fraction differences are absolute, using fibroblast-restricted summaries and the unchanged threshold. The split-half column lists the 97.5th percentile of absolute mean differences from 250 random disjoint fibroblast partitions within each specimen, in pair order. The empirical CDF supremum compares the two full loading distributions. Cell partitions describe conditional finite-cell variability; they are not biological replicates, confidence intervals, a patient-level null, an intraclass correlation or clinical reliability estimates.}\label{tab:p1repeats}\\
\toprule
Specimens & Group & Mean difference & Split-half limits & High-fraction diff. & CDF supremum \\
\midrule\endfirsthead
\multicolumn{6}{l}{\tablename\ \thetable\ (continued)}\\
\toprule
Specimens & Group & Mean difference & Split-half limits & High-fraction diff. & CDF supremum \\
\midrule\endhead
\bottomrule\endfoot
G23 / G4 & Healed & 0.0096 & 0.0309 / 0.0445 & 0.0025 & 0.0576 \\*
G33 / G34 & Not healed & 0.0042 & 0.0017 / 0.0041 & 0.0053 & 0.1251 \\
G7 / G8 & Healed & 0.2583 & 0.0553 / 0.0185 & 0.5458 & 0.5789 \\
G31 / G32 & Non-diabetic & 0.0002 & 0.0007 / 0.0009 & 0.0000 & 0.0337 \\*
G10 / G24 & Non-diabetic & 0.0019 & 0.0010 / 0.0013 & 0.0000 & 0.1865 \\
\end{longtable}
}

{\small
\begin{longtable}{@{}p{37mm}p{18mm}p{18mm}p{20mm}p{33mm}p{36mm}@{}}
\caption{Revised same-model held-pair discrimination in 11 mapped DFU patients (7 healing, 4 non-healing). Each of the 28 folds holds one patient from each arm out, fits on the other nine and scores both with one unchanged rule. Wins favour the healing patient; an exact tie earns one-half credit. Discrimination is total credit divided by 28, not a pooled-fold raw-coordinate AUC. Topic and module orientation and PCA/NMF fitting and component selection use training patients. NMF basis normalization and training-SD-standardized component selection remove positive scaling ambiguity. The uniform NMF cap was raised from 1,000 to 10,000 after training-convergence warnings; no final fit warned. All pairs share 11 patients, and discovery expression informed the fixed panel and encoder. No confidence intervals, p values, independent validation or representation ranking are claimed.}\label{tab:p1pairauc}\\
\toprule
Readout & Wins & Ties & Losses & Pair credit & Discrimination \\
\midrule\endfirsthead
\multicolumn{6}{l}{\tablename\ \thetable\ (continued)}\\
\toprule
Readout & Wins & Ties & Losses & Pair credit & Discrimination \\
\midrule\endhead
\bottomrule\endfoot
Frozen topic & 20 & 0 & 8 & 20.0 / 28 & 0.7143 \\*
Module & 22 & 0 & 6 & 22.0 / 28 & 0.7857 \\
PCA & 14 & 0 & 14 & 14.0 / 28 & 0.5000 \\*
NMF & 5 & 3 & 20 & 6.5 / 28 & 0.2321 \\
\end{longtable}
}

{\small
\begin{longtable}{@{}p{25mm}p{27mm}p{25mm}p{31mm}p{30mm}p{27mm}@{}}
\caption{Full released-count depth audit of all 81,602 pre-analysis-QC columns from 25 foot-skin specimens. Baseline reproduces 76,461 retained cells and 19,410 fibroblast labels. Each nonzero count is independently binomial-thinned at the stated probability; each rate starts from original counts and uses all three fixed computational seeds. QC is at least 200 detected genes and mitochondrial fraction at most 0.20. Marker means/SDs, forced-winner gate, encoder and high-state cut remain frozen. Lost and gained cells are relative to baseline QC, not a net count. Lineage switches divide changed labels by cells surviving both QC passes. Zero-marker cells pass QC but detect none of the annotation markers; their forced labels are not independently validated. Counts are cells, not independent subjects; seeds provide no biological replication.}\label{tab:p1fulldepth}\\
\toprule
Count retention & Seed & QC cells & Lost / gained & Lineage switches & Zero-marker cells \\
\midrule\endfirsthead
\multicolumn{6}{l}{\tablename\ \thetable\ (continued)}\\
\toprule
Count retention & Seed & QC cells & Lost / gained & Lineage switches & Zero-marker cells \\
\midrule\endhead
\bottomrule\endfoot
100\% & Baseline & 76461 & 0 / 0 & 0.00\% & 92 \\*
75\% & 20260925 & 76375 & 127 / 41 & 3.18\% & 165 \\
75\% & 20260926 & 76387 & 113 / 39 & 3.30\% & 169 \\
75\% & 20260927 & 76393 & 117 / 49 & 3.29\% & 168 \\
50\% & 20260925 & 75738 & 793 / 70 & 6.23\% & 384 \\
50\% & 20260926 & 75726 & 799 / 64 & 6.25\% & 363 \\
50\% & 20260927 & 75724 & 810 / 73 & 6.34\% & 386 \\
25\% & 20260925 & 72145 & 4419 / 103 & 10.29\% & 770 \\
25\% & 20260926 & 72197 & 4367 / 103 & 10.48\% & 842 \\
25\% & 20260927 & 72145 & 4404 / 88 & 10.45\% & 815 \\
10\% & 20260925 & 61397 & 15175 / 111 & 15.36\% & 1429 \\
10\% & 20260926 & 61366 & 15210 / 115 & 15.40\% & 1462 \\*
10\% & 20260927 & 61395 & 15178 / 112 & 15.31\% & 1410 \\
\end{longtable}
}

{\small
\begin{longtable}{@{}p{25mm}p{48mm}p{26mm}p{35mm}p{33mm}@{}}
\caption{Patient contrasts after rerunning QC and the frozen-reference gate on each depth perturbation. Differences are healed minus non-healed, with seven and four author-mapped DFU patients. Repeated specimens are pooled by the appropriate all-cell or fibroblast count. Fibroblast composition uses all current QC survivors; high state uses the unchanged discovery cut. The exact two-sided probability enumerates all 330 patient-label allocations, including ties, without add-one correction. Baseline has one run; other rows give the minimum and maximum over three fixed thinning seeds. These are computational ranges, not confidence intervals, multiplicity-adjusted discoveries, equivalence bounds or biological replication. Zero-panel cells remain explicitly missing in conditional branches rather than receiving scores from an all-zero vector.}\label{tab:p1depthoutcomes}\\
\toprule
Count retention & Readout & Estimable runs & Difference range & Exact p range \\
\midrule\endfirsthead
\multicolumn{5}{l}{\tablename\ \thetable\ (continued)}\\
\toprule
Count retention & Readout & Estimable runs & Difference range & Exact p range \\
\midrule\endhead
\bottomrule\endfoot
100\% & All-cell mean & 1 / 1 & 0.0293 to 0.0293 & 0.2727 to 0.2727 \\*
100\% & Fibroblast mean & 1 / 1 & 0.0892 to 0.0892 & 0.3121 to 0.3121 \\
100\% & High-state fraction & 1 / 1 & 0.1514 to 0.1514 & 0.3697 to 0.3697 \\
100\% & Fibroblast composition & 1 / 1 & 0.0708 to 0.0708 & 0.4697 to 0.4697 \\
75\% & All-cell mean & 3 / 3 & 0.0291 to 0.0292 & 0.2758 to 0.2818 \\
75\% & Fibroblast mean & 3 / 3 & 0.0876 to 0.0899 & 0.2970 to 0.3182 \\
75\% & High-state fraction & 3 / 3 & 0.1451 to 0.1518 & 0.3545 to 0.3788 \\
75\% & Fibroblast composition & 3 / 3 & 0.0687 to 0.0706 & 0.4727 to 0.4788 \\
50\% & All-cell mean & 3 / 3 & 0.0289 to 0.0291 & 0.2818 to 0.2848 \\
50\% & Fibroblast mean & 3 / 3 & 0.0880 to 0.0913 & 0.2848 to 0.3152 \\
50\% & High-state fraction & 3 / 3 & 0.1485 to 0.1620 & 0.3364 to 0.3727 \\
50\% & Fibroblast composition & 3 / 3 & 0.0683 to 0.0698 & 0.4758 to 0.4848 \\
25\% & All-cell mean & 3 / 3 & 0.0292 to 0.0293 & 0.2788 to 0.2879 \\
25\% & Fibroblast mean & 3 / 3 & 0.0883 to 0.0902 & 0.2758 to 0.2758 \\
25\% & High-state fraction & 3 / 3 & 0.1526 to 0.1634 & 0.3303 to 0.3424 \\
25\% & Fibroblast composition & 3 / 3 & 0.0737 to 0.0759 & 0.4455 to 0.4727 \\
10\% & All-cell mean & 3 / 3 & 0.0288 to 0.0296 & 0.3303 to 0.3455 \\
10\% & Fibroblast mean & 3 / 3 & 0.0879 to 0.0918 & 0.1970 to 0.2212 \\
10\% & High-state fraction & 3 / 3 & 0.1645 to 0.1687 & 0.2515 to 0.3000 \\*
10\% & Fibroblast composition & 3 / 3 & 0.0845 to 0.0897 & 0.3636 to 0.4091 \\
\end{longtable}
}
